% =====================================================================
% Table 1 --- the two-axis core of the closure matrix, block-tagged.
% 2026-10-27: merged with the former overview table (archived); renders as Table 1.
% ---------------------------------------------------------------------
% Revised 2026-10-18 (ADR-0033): the financing conditions name the
% balancing variable. tau is the wage-income tax rate that finances the
% baseline government purchases g^G in every closure (kernel tau_rate);
% under F2 the programme g is added to that base. F is the net
% external transfer; B_gov = sum p g is the programme booked externally (F3).
% Block tags \textbf{labor market:} / \textbf{financing:} sit on their own
% line, so the two "financing:" labels align by construction.
% =====================================================================
\begin{table}[htbp]
\centering
\footnotesize
\setlength{\tabcolsep}{4pt}
\setlength{\parindent}{0pt}
\setlength{\parskip}{0pt}
\renewcommand{\arraystretch}{1.3}
\newcommand{\tblk}[1]{\textbf{#1:}}
\newcommand{\tcond}[1]{\par\hspace{1em}#1}
\newcommand{\outcome}[1]{{\centering\itshape #1\par}}
\begin{tabular}{@{}>{\raggedright\arraybackslash}p{2.2cm}
                  >{\raggedright\arraybackslash}p{6.2cm}
                  >{\raggedright\arraybackslash}p{6.2cm}@{}}
\toprule
 & \multicolumn{2}{c}{\textbf{Labor-market closure}} \\
\cmidrule(lr){2-3}
\textbf{Financing closure}
& \textbf{Standard CGE}\newline \scriptsize one flexible wage, labor mobile, full employment
& \textbf{Standard Leontief}\newline \scriptsize wage pinned, fixed coefficients, output absorbs \\
\midrule
\textbf{F1}\newline \scriptsize repurposing of consumption
&             {\outcome{pure reallocation}}
  \vspace{5pt}
\tblk{labor market}\tcond{$L_{i} = L_{i}^{\mathrm{cm}}(p, y_{i}, w)$, \quad $\textstyle\sum_{i} L_{i} = \bar L$}\par
  \vspace{6pt}
  \tblk{financing}\tcond{$\tau\, w\sum_{i} L_{i} = \textstyle\sum_{i} p_{i} g^{G}_{i}$, \quad $\tilde\beta_{i}$ tilted to program sectors, \quad $B_{\mathrm{gov}} = 0$}
& {\outcome{quantity adjustment}}
  \vspace{5pt}
\tblk{labor market}\tcond{$w = 1$, \quad $L_{i} = L_{i}^{\mathrm{cm}}(p, y_{i}, w)$, \quad $(\theta, \varepsilon, \sigma) \to 0^{+}$}\par
  \vspace{6pt}
  \tblk{financing}\tcond{$\tau\, w\sum_{i} L_{i} = \textstyle\sum_{i} p_{i} g^{G}_{i}$, \quad $\tilde\beta_{i}$ tilted to program sectors, \quad $B_{\mathrm{gov}} = 0$} \\
\midrule
\textbf{F2}\newline \scriptsize refinancing via taxes
& {\outcome{crowding out}}
  \vspace{5pt}
\tblk{labor market}\tcond{$L_{i} = L_{i}^{\mathrm{cm}}(p, y_{i}, w)$, \quad $\textstyle\sum_{i} L_{i} = \bar L$}\par
  \vspace{6pt}
  \tblk{financing}\tcond{$\tau\, w\sum_{i} L_{i} = \textstyle\sum_{i} p_{i}(g^{G}_{i} + g_{i})$, \quad $B_{\mathrm{gov}} = 0$}
& {\outcome{crowding out with a Haavelmo top-up}}
  \vspace{5pt}
\tblk{labor market}\tcond{$w = 1$, \quad $L_{i} = L_{i}^{\mathrm{cm}}(p, y_{i}, w)$, \quad $(\theta, \varepsilon, \sigma) \to 0^{+}$}\par
  \vspace{6pt}
  \tblk{financing}\tcond{$\tau\, w\sum_{i} L_{i} = \textstyle\sum_{i} p_{i}(g^{G}_{i} + g_{i})$, \quad $B_{\mathrm{gov}} = 0$} \\
\midrule
\textbf{F3}\newline \scriptsize external finance
& {\outcome{crowding out}}
  \vspace{5pt}
\tblk{labor market}\tcond{$L_{i} = L_{i}^{\mathrm{cm}}(p, y_{i}, w)$, \quad $\textstyle\sum_{i} L_{i} = \bar L$}\par
  \vspace{6pt}
  \tblk{financing}\tcond{$\tau\, w\sum_{i} L_{i} = \textstyle\sum_{i} p_{i} g^{G}_{i}$, \quad $B_{\mathrm{gov}} = \textstyle\sum_{i} p_{i} g_{i}$}
& {\outcome{output expansion and external debt}}
  \vspace{5pt}
\tblk{labor market}\tcond{$w = 1$, \quad $L_{i} = L_{i}^{\mathrm{cm}}(p, y_{i}, w)$, \quad $(\theta, \varepsilon, \sigma) \to 0^{+}$}\par
  \vspace{6pt}
  \tblk{financing}\tcond{$\tau\, w\sum_{i} L_{i} = \textstyle\sum_{i} p_{i} g^{G}_{i}$, \quad $B_{\mathrm{gov}} = \textstyle\sum_{i} p_{i} g_{i}$} \\
\bottomrule
\end{tabular}
\caption{The two-axis core of the closure matrix: the three financing closures
(rows) crossed with the two labor-market poles of the introduction (columns).
Each cell states the labor-market condition (the column) and the financing
condition (the row), and the outcome line states the qualitative response to
the demand-only impulse; the program cost $\sum_i p_i g_i$ is valued at each
cell's equilibrium prices. In the conditions, $\tau$ is the wage-income tax rate financing the
baseline purchases $g^{G}$ (a block unchanged by the program), $g$ is the
program bundle, $L_{i}^{\mathrm{cm}}$ is cost-minimizing labor demand, and
$\tilde\beta_{i}$ are expenditure shares tilted toward the program sectors
under F1.}
\label{tab:closure-conditions}
\end{table}
